\scanpage{86}
% archive sheet number: 85
% Face upside down in the scan. Long oblique cancellation strokes cover the typescript.
\deleted{\deleted{\unsure{Notons}} Notons en passant le résultat suivant, qui servira dans un exposé ultérieur:}

\deleted{\uline{Corollaire} \handwritten{4.6.} Soient $X$ un $S$-préschéma plat, $U$ une partie ouverte de $X$, on suppose que pour tout $s\handwritten{\in} S$, $U_s$ est schématiquement dense dans $X_s$.
Supposons de plus que $X$ est localement noethérien, \uline{ou} $X$ localement de présentation finie sur $S$ et $U$ rétrocompact dans $X$ (ou ce qui revient au même, $U$ de présentation finie sur $X$; il suffit par exemple $X$ et $U$ de présentation finie sur $S$).
Alors $U$ est schématiquement dense dans $X$.}

\deleted{Le cas $X$ localement noethérien n'est mis que pour mémoire, il est contenu dans \handwritten{4.3.} Dans le deuxième cas envisagé, on peut évidemment
supposer $S$ et $X$ affine, alors $X$ et $U$ sont de \unsure{présentattion|présentation} finie sur $S=\operatorname{Spec}(A)$. \inserted{$A$ est limite inductive de ses sous-$\underline{Z}$-algèbres $A_i$ de type fini.} Le ``procédé bréveté'' déj\unsure{è|e|à} utilisé (explicité en long et en
large dans EGA IV 8) montre qu'il existe un \deleted{schéma noethérien affine} \deleted{\unsure{et un anneau|et un|sous-anneau}} 
$i$, un schéma affine $X_i$ sur $S_i=\operatorname{Spec}(A_i)$ et un ouvert
$U_i$ dans $X_i$, dont $X,U$ se déduisent par changement de base \deleted{\unsure{et|de|si}} $S\quad S_i$.
\deleted{Pour tout $j\quad i$,} soit $E_{\deleted{\unsure{j|i}}\handwritten{i}}$ la partie de $S_{\handwritten{i}}$ formée des $s\quad S_{\handwritten{i}}$ tels que $U_{\handwritten{i}\,s}$ soit schématiquement dense dans $X_{\handwritten{i}\,s}$. Cette condition signifie aussi
 que $U_{\handwritten{i}\,s}$ contient \deleted{\unsure{les|le}} l'ensemble $\operatorname{Ass}\,\underline{O}_{X_{\handwritten{i}\,s}}$ des points ``associés'' au faisceau structural de $X_{\handwritten{i}\,s}$, et les résultats de EGA IV 9 montrent
alors que cette condition est de nature constructible, i.e. que $E_{\handwritten{i}}$ est une partie constructible de $S_{\handwritten{i}}$. D'ailleurs, la question envisagée sur
les fibres est invariante par extension du corps de base, ce qui prouve \deleted{\unsure{que pour $k,j,$|que pour $k$|que pour tout $j$}\illegible{2}} \inserted{en vertu de l'hypothèse faite que} 
l'image inverse \inserted{$\handwritten{\unsure{E|=}}$} de $E_{\handwritten{\unsure{i|j}}}$ par $S_{\deleted{\unsure{k|j}}}\handwritten{\longrightarrow} S_{\handwritten{\unsure{i|j}}}$ \deleted{\unsure{par que|part que|part}} est
$S$, et entraine par EGA IV 8 qu'il existe un $j\handwritten{\geq}i$ tel que l'image
inverse \deleted{\unsure{de|le}} $E_j$ de $E_i$ dans $S_j$ soit $S_j$, i.e. que pour tout $s\handwritten{\in}S_j$, $U_{j\,s}$ est schématiquement dense dans $X_{j\,s}$. 
En vertu de 4.4. appliqué à $(\unsure{S|E}_j,X_j,U_j$ et au\deleted{\unsure{x}} changement de base $S\handwritten{\to} S_j$, il s'ensuit que $\unsure{U|X}$ est schématiquement dense dans $X$.
On notera d'ailleurs qu'on n'utilise ici 4.4. que dans le}
% Last line ends at the edge after « le ».
